\documentclass[conference]{IEEEtran}

\usepackage{amsmath}
\usepackage{amssymb}
\usepackage{booktabs}
\usepackage{graphicx}
\usepackage{cite}
\usepackage{url}
\usepackage{microtype}
\usepackage{array}

\graphicspath{{../}{./}}

\begin{document}

\title{Acoustic Gunshot Detection using Edge AI, TinyML and a Wireless Sensor Network:\\ An Empirical Journey from Hand-Crafted Features to On-Device Binary Inference}

\author{
    \IEEEauthorblockN{Jeevant Sharma, Deepesh Dangi, Abhishek Yadav, Bhupendra Kumar}
    \IEEEauthorblockA{\textit{Department of Computer Science and Engineering}\\
    \textit{Indian Institute of Information Technology, Surat}\\
    Under the guidance of Dr.~Kaustubh Dhondge}
}

\maketitle

\begin{abstract}
Gun violence in a densely populated city demands a response measured in seconds, not minutes, and yet most
discharges are never reported and the ones reported almost always give an inaccurate location. We describe
an inexpensive, privacy-preserving Wireless Sensor Network (WSN) in which every node runs an on-device
classifier; raw audio never leaves the device and only a short binary alert is transmitted, after which
multiple anchor nodes let a hub localise the event using Time Difference Of Arrival (TDOA). This manuscript
is written as an honest chronological field log rather than a single cherry-picked model: we trace the path
from raw data cleaning, through 250\,ms/500\,ms/1\,s trimming, class-imbalance studies from 1:1 to 1:50,
supervised baselines (linear regression, SVM, Random Forest) that scored 99--100\% on closed splits and then
collapsed to near-constant false alarms on a live microphone, into the Edge Impulse TinyML experiments (MFCC
and MFE spectrogram, 1D and 2D CNN, anomaly detection), then into the native deep-CNN progression
(1D $\rightarrow$ 2D $\rightarrow$ 3D $\rightarrow$ enhanced dual-head), and finally into a bioacoustic
Per-Channel Energy Normalization (PCEN) $+$ Bidirectional GRU recurrent network. On a balanced, unseen
300-clip in-the-wild benchmark (100 real gunshots, 100 ambient, 100 hard imposters) the lightweight Enhanced
2D CNN reached sub-millisecond inference ($0.09$\,ms, $0.05\%$ of the $187.5$\,ms hop budget) with up to
$99\%$ ambient specificity, the Robust CRNN-PCEN showed zero recall degradation under synthetic microphone
distortion, and a multi-model ensemble reached $77\%$ gunshot recall by combining complementary
representations with physics-based acoustic transient rejection. We close with the deployment roadmap---%
sleep-scheduled nodes, a 1-bit (gunshot/no-gunshot) binary alert that keeps radio energy near zero,
sink-side spatial aggregation, and a live dashboard---and with the physical boundary of a single microphone,
which is mortar fireworks.
\end{abstract}

\begin{IEEEkeywords}
Acoustic gunshot detection, Edge AI, TinyML, PCEN, CRNN, Wireless Sensor Network, TDOA, quantization,
binary inference.
\end{IEEEkeywords}

% ==============================================================================
\section{Introduction}
\label{sec:intro}
% ==============================================================================

\subsection{Motivation and Societal Context}
In the US and Europe people who carry guns can have a licence, and even in India this is true for the big
cities, while in some other places people do it illegally. For the initiative part---to start the well-being
of a typical city---we must, as far as we can, save lives and be able to let, or at least decrease, the abuse
that people do in the streets. Since we are in a digital age, we can have a different, much better set of
tools and devices that are able to capture the information of those gun sounds, along with the heavy
processing needed to send that data from widely spread microcontrollers, and then speed it up to the forest
and other national sanctuaries to preserve wildlife as well as the lives of humans.

The police are, typically, a very powerful and helpful way to make far ahead of the typical way of living in
the real world along with technology. Even with the technologies that are present---we applied a chainsaw
detection, and other services also provide it---the concept of a smart city aligns with our goal. If people
are not co-operating, the system does this job. If in India, due to heavy crowd and a more abusive community
in some places, it is impossible for now, then at some states it can be started with colleges or school
institutions, so that it truly works the way we expect. A small idea of this problem, and the conclusion of
this introduction, is this: the gunshots that are \textit{for the society are a typical example of a bad
will}---but with technology, not using humans, we can detect these fires, we can easily let people know where
the event happens, and the person who gets injured can get medical support. Our law ideology takes the shape
of those ways that are like catching criminals, plus taking care of the loss caused by the weapon to society,
with the help of tags along with IoT.

\subsection{What We Set Out to Build}
We like to do the typical learning and run the inference on the Arduino BLE. But with the work and the
real-world problem, we started the way things natively go---a simple run on a PC/laptop after training, then
on the PC, then shifted to the Arduino with models converted from \texttt{.h5} to \texttt{.tflite}. Our goal
is to create and run the inference on the Arduino BLE. In this paper we worked up to a certain point around
the typical deployment, and typically we ask it to do that job through inference, and we need to send audio
to it just as a say of $1$ and $0$ in that case. This was also tried on the Raspberry Pi, on the virtual one,
with far fewer settings, and we record these as a future aspect we are still working on, while our thresholds
are still where we are working.

\subsection{Contributions}
\begin{enumerate}
    \item A complete engineering log of an acoustic gunshot detector built end-to-end on cheap hardware,
    including every dead end (250\,ms truncation, classical-ML collapse, the ``Dying ReLU'', the MFCC
    voice-trap, and the INT8 quantization failure of raw-waveform models).
    \item A systematic Edge Impulse TinyML comparison of MFCC vs.\ MFE representations and 1D vs.\ 2D
    convolution, with on-device RAM/Flash/latency footprints.
    \item A bioacoustic PCEN $+$ Bi-GRU front-end proven, by measurement, to be invariant to synthetic
    microphone domain shift.
    \item A balanced 300-clip in-the-wild evaluation across five INT8 architectures under two windowing
    paradigms, plus a multi-model ensemble reaching $77\%$ recall.
    \item An edge-native WSN design---sleep scheduling, a 1-bit binary alert, and sink-side spatial
    consensus---making dense, private, battery-friendly deployment feasible.
\end{enumerate}

% ==============================================================================
\section{Related Work}
\label{sec:related}
% ==============================================================================
\textbf{Commercial acoustic sensing.} Systems such as ShotSpotter proved that acoustic detection works in a
city, but they are proprietary, cloud-tethered and expensive; dense deployment is unaffordable, and streaming
raw audio off-node raises real privacy and capital-cost concerns.

\textbf{Classical acoustic detection.} Much of the literature uses hand-crafted features---Mel-Frequency
Cepstral Coefficients (MFCC), spectral centroid, spectral rolloff, zero-crossing rate---with Support Vector
Machines or Random Forests. These report excellent closed-set accuracy but, as we confirm in
Section~\ref{sec:classical}, average the signal in time and destroy the micro-second shockwave attack that
physically defines a gunshot.

\textbf{Deep spectrogram models.} Convolutional networks over time-frequency representations (log-Mel, MFE)
dominate modern audio tagging and suit TinyML, but are fragile to microphone domain shift and to unnormalized
decibel inputs.

\textbf{Bioacoustic / DCASE lineage.} Our main inspiration is rainforest bioacoustics and illegal-chainsaw
detection: Per-Channel Energy Normalization (PCEN) with a Bidirectional Gated Recurrent Unit. PCEN acts as a
temporal automatic gain control that flattens stationary background noise and absorbs transducer coloration.

\textbf{TinyML and edge deployment.} The MCUNet / MobileNet philosophy---small receptive fields, aggressive
INT8 quantization, a few hundred KB---guides our on-device targets. To the best of our knowledge, combining
(i) a PCEN-CRNN front-end, (ii) an empirically verified domain-shift invariance study, and (iii) a 1-bit
binary-alert WSN protocol in a single open deployment is a unique combination (see Section~\ref{sec:unique}).

% ==============================================================================
\section{Implementation and Experimentation}
\label{sec:method}
% ==============================================================================
\textit{Note: the experimentation is described here, and all measurements from every stage are collected in
Section~\ref{sec:results}, exactly as the handwritten log intended (implementation first, results second).}

\subsection{Data Collection, Cleaning and Manual Auditing}
We started with the cleaning of data, where all data was passed through the pipeline. Cleaning used
\textbf{librosa} and many different ways to look along, and was a major task in ML. We took data from many
sources, then trimmed it along \textbf{250\,ms, 500\,ms and 1\,s} segments, and removed useless parts. The
dataset contained roughly \textbf{80K files}, so we checked about \textbf{1K--3K--5K} files by manually
listening in cycles until we were satisfied that the data was cleaned well. Sources were a calibrated
multi-firearm benchmark, field firearm recordings (handguns, rifles, shotguns) and long-duration negatives
(urban traffic, rainfall, construction, speech, dogs, room tone).

\subsection{The 250\,ms Mistake, Onset Detection and Quality Gates}
Our initial trimmers used a fixed \textbf{250\,ms} window around the energy peak, which proved to be a
critical flaw. A real discharge has the supersonic \textbf{shockwave (N-wave)} ($\approx$~1--5\,ms pressure
rise) and the \textbf{muzzle blast} with a reverberation tail of $\approx$~200--600\,ms. A 250\,ms window cuts
the tail, so a gunshot starts to look like a door slam or a clap. We therefore standardised to
\textbf{750\,ms} ($N=16{,}537$ samples at $f_s=22{,}050$\,Hz) and built a 3-way onset detector:
\begin{itemize}
    \item \textbf{Spectral flux} (half-wave rectified frame-to-frame surge):
    \begin{equation}
        O(m) = \sum_{k} \max\left(0, |X(k,m)| - |X(k,m-1)|\right)
    \end{equation}
    \item \textbf{Raw amplitude peak}: $n_{\text{peak}} = \arg\max_n |y[n]|$.
    \item \textbf{Smoothed envelope}: moving average of $|y|$ with kernel $K = 0.005\,f_s$.
\end{itemize}
Every candidate slice passes three physical quality gates:
\begin{equation}
    \text{Crest Factor} = \frac{\max(|y|)}{\text{RMS} + \epsilon}
\end{equation}
\begin{equation}
    \text{Prominence} = \frac{\max(|y|)}{\text{Median}(|y|) + \epsilon}
\end{equation}
\begin{equation}
    \text{Attack Ratio} = \frac{\sum_{n=0}^{N/2-1} y[n]^2}{\sum_{n=N/2}^{N-1} y[n]^2 + \epsilon}
\end{equation}
Negatives with a high Attack Ratio and extreme Crest Factor were purged, so the ambient class did not
secretly contain impulses.

\subsection{Class Imbalance: From 1:1 to 1:10 to 1:50}
We tried the models with different dataset splits---starting from $\approx$~1:1, going to \textbf{1:10}, and
finally understanding that total gunshot capability can realistically sit anywhere between \textbf{1:10 and
1:50}. Models trained strictly on 1:1 were over-sensitive and produced heavy false alarms during continuous
monitoring; re-weighting the loss and enriching the negative corpus (sirens, traffic, wind, thunder, dog
barks, doors) was required to steady the boundary.

\subsection{Why Classical Supervised Learning First---and Why It Failed}
\label{sec:classical}
We started with the basic question of whether a simple linear decision boundary could separate the classes
at all, and used it to build a live system to get some thresholds and see which model can be the
lowest-latency one. We then saw \textit{why and how it falls}, and moved to \textbf{SVM} and \textbf{Random
Forest}, trained on 13-band MFCC features with balanced class weights. On a closed split SVM (RBF) reached
$\approx 0.953$ accuracy and Random Forest $\approx 0.99$; but deployed to a live microphone both collapsed
into continuous false alarms on percussive sounds (claps, keys, doors, brakes). The reason is physical:
computing MFCCs averages spectral frames in time, which smears the sub-3\,ms shockwave attack into a blurred
frequency envelope --- discarding the one signature that separates firearms from everyday percussion.

\subsection{The Edge Impulse TinyML Exploration}
Before the native CNNs, we moved to \textbf{Edge Impulse}, preprocessing CNN spectrograms and testing
several field models. We saw a lot better accuracy for 2D models but \textbf{very bad results on live
audio}, and that problem stayed part of our flow. Five experiments were run on the ARM Cortex-M4F target
(64\,MHz, 256\,KB SRAM):
\begin{itemize}
    \item \textbf{Exp.~1 (MFCC $+$ 1D CNN, the ``voice-trap'')}: 250\,ms window, frame 0.025\,s, stride
    0.02\,s, 32 Mel filters, 13 MFCCs, FFT 512; validation \textbf{99.4\%}, INT8 1\,ms, RAM 11.9\,KB, Flash
    45.4\,KB. Live: \textbf{33/47 windows false-triggered ($\sim$70\%)} on speech/ambient.
    \item \textbf{Exp.~2 (MFCC $+$ K-means anomaly)}: DSP 120\,ms; $47/48$ windows flagged as anomalies but
    the classifier still false-fired --- the MFCC representation lacked temporal discriminability.
    \item \textbf{Exp.~3 (MFE spectrogram $+$ 1D conv)}: FFT 32, 3{,}383 features; validation
    \textbf{98.8\%}, INT8 12\,ms, Flash 47.1\,KB. Live: 24/48 false alarms ($50\%$); the voice-trap broke,
    but 1D convolution over unrolled bins cannot capture 2D structure.
    \item \textbf{Exp.~4 (MFE spectrogram $+$ 2D CNN, the turning point)}: validation \textbf{99.6\%},
    ROC-AUC 1.00, INT8 167\,ms, RAM 30.2\,KB, Flash 52.2\,KB; heavy wind $\rightarrow$ 22 rejections vs 7
    false alarms ($75.8\%$), thunder $\rightarrow$ 20 vs 5 ($80.0\%$), true 12-gauge shot $\rightarrow$
    probability $0.06 \to >0.94$ at the blast onset.
    \item \textbf{Exp.~5 (MFE spectrogram $+$ anomaly)}: validation \textbf{99.7\%}, a second rejection gate
    for non-firearm high-amplitude events.
\end{itemize}
At the end we created a report, and the time was up for these pre-existing models --- it was now on us to use
the CNN models in our own way. So we looked for CNN architectures starting with a \textbf{1D CNN} and moved
on: \textbf{2D CNN $\rightarrow$ 3D CNN $\rightarrow$ enhanced}.

\subsection{Native Deep CNN Progression and Hardware Conditioning}
\textbf{Baseline 1D raw-waveform CNN (92{,}513 params).} The first layer is 80-sample wide
($K=80$, stride 4, 32 filters), spanning
\begin{equation}
    \Delta t = \frac{80}{22{,}050\,\text{Hz}} \approx 3.63\,\text{ms},
\end{equation}
a learnable bank of 32 FIR wavelets matched to N-wave rise time. \textit{Failure mode:} 99.7\% on clean
held-out files but recall collapsed to \textbf{30.0\%} across different microphones and rooms.

\textbf{Baseline 2D Mel CNN (249{,}985 params).} 64 Mel bins $\times$ 130 frames (FFT 512, hop 128), four
Conv2D stages (32$\to$64$\to$128$\to$128). \textit{Failure mode (Dying ReLU collapse):} decibel values in
$[-80,0]$ drove ReLU activations to zero and killed all gradients (0\% recall, 679 false negatives);
normalizing to the unit range,
\begin{equation}
    S_{\text{norm}} = \frac{S_{\text{dB}} + 80}{80},
\end{equation}
restored convergence to 99.86\% ROC-AUC.

\textbf{Enhanced dual-head models.} MixUp
($\tilde{x} = \lambda x_i + (1-\lambda)x_j$, $\lambda \sim \text{Beta}(0.3,0.3)$), SpecAugment time masking,
a second ``acoustic anomaly'' head, and a class-weighted $F_2$ loss that penalises false negatives far more
than false alarms ($C_1 \approx 6.80$, $C_0 \approx 0.54$).

\textbf{Hardware reality.} Low-cost microphones (INMP441, MAX9814) introduce a DC rail offset (baseline
$\approx +0.05/-0.08$), thermal ADC hiss, and diaphragm clipping above $\approx 120$\,dB SPL. We strip DC by
mean subtraction and gate peak normalization by energy:
\begin{equation}
    y_{\text{clean}}[n] = y[n] - \frac{1}{N}\sum_{i=0}^{N-1} y[i],
\end{equation}
\begin{equation}
    y_{\text{norm}} = \begin{cases}
        \dfrac{y}{\max(|y|)} & \text{RMS} \ge 0.005 \\[6pt]
        y & \text{RMS} < 0.005
    \end{cases}
\end{equation}
We also frame with a 750\,ms ring buffer at 75\% overlap (187.5\,ms hop) so no transient is bisected.

\subsection{PCEN $+$ CRNN: The Bioacoustic Adaptation}
Giving some time to these models, we got a research paper on \textbf{PCEN (Per-Channel Energy
Normalization)}, and we tried to work it out in a typical way. PCEN needs a different kind of cleaning of the
file before you touch it. The front-end is STFT (FFT 512, hop 128), a \textbf{Slaney} Mel filterbank (64
bands, bandwidth-normalized, $\text{FilterWeight}_m(f) = \text{Triangle}_m(f)\cdot 2/(f_{m+2}-f_m)$), a
$2^{31}$ scale, a causal IIR envelope
\begin{equation}
    M[f,t] = (1-b)\,M[f,t-1] + b\,S[f,t], \qquad b \approx 0.2068,
\end{equation}
and the PCEN non-linearity
\begin{equation}
    P[f,t] = \left(\frac{S[f,t]}{(\epsilon + M[f,t])^{\alpha}} + \delta\right)^{r} - \delta^{r},
\end{equation}
with $\alpha = 0.98$, $\delta = 2.0$, $r = 0.5$, $\epsilon = 10^{-6}$, then
$X = (P - 0.439023)/0.187652$. Under stationary noise the steady-state gain converges to
\begin{equation}
    \lim_{t\to\infty} E(t) \approx S_0^{1-\alpha} = S_0^{0.02} \approx 1.0,
\end{equation}
and under a microphone gain $g$ it scales as $g^{0.04} \approx 1.0$ (transducer gain invariance). Features
are sequenced into 32 frames ($D=1024$) and read by a Bidirectional GRU ($2\times64$): the forward pass
tracks the attack, the backward pass verifies the reverberation decay.

\subsection{From \texttt{.h5} to \texttt{.tflite}: Quantization and the FlexOp Crisis}
Converting the CRNN \texttt{.h5} to TFLite produced a \texttt{FlexTensorListReserve} operation that a light
runtime cannot prepare. We fixed this by \textbf{statically unrolling the Bi-GRU across 32 fixed time steps}
and applying INT8 post-training quantization, yielding a fully native LiteRT model. We also fixed three
preprocessing mismatches found on in-the-wild audio: missing z-score normalization for the Baseline 2D CNN
(0\% $\to$ 49.2\% recall), a wrong filterbank for the Enhanced 2D CNN (27.1\% $\to$ 44.1\% recall, latency
2.5\,ms $\to$ 0.28\,ms), and a non-Slaney SciPy PCEN whose correction lifted real-shot confidence from
$1.6\% \to 95.6\%$ (matched to \texttt{librosa.pcen} to within $8.45\times10^{-8}$).

\subsection{Evaluation Protocol}
Final evaluation used a balanced, unseen \textbf{300-clip in-the-wild set} (100 real gunshots, 100 ambient,
100 hard imposters) from FSD50K, with zero training overlap. Two paradigms were compared: \textbf{peak-centered}
(one 750\,ms slice at the loudest sample) and \textbf{sliding-window} (750\,ms, 75\% overlap). Hardware
profiling used a Raspberry Pi 4 class target (a VirtualBox VM throttled to 45--50\% CPU after physical Pi boot
failures) with the LiteRT INT8 runtime, 100 iterations per model.

% ==============================================================================
\section{Results}
\label{sec:results}
% ==============================================================================

\begin{table}[t]
\centering
\caption{Edge Impulse experiments on Cortex-M4F (INT8).}
\label{tab:ei}
\small
\begin{tabular}{p{0.6cm}p{2.4cm}p{1.5cm}c c c c}
\toprule
\textbf{Exp} & \textbf{Representation} & \textbf{Classifier} & \textbf{Val.} & \textbf{Lat.} & \textbf{RAM} & \textbf{Flash} \\
\midrule
1 & 13 MFCCs & 1D CNN & 99.4\% & 1\,ms & 11.9\,KB & 45.4\,KB \\
2 & MFCC + K-means & 1D + Anom. & 99.4\% & 2\,ms & 11.9\,KB & 45.4\,KB \\
3 & MFE (3{,}383) & 1D CNN & 98.8\% & 12\,ms & 18.1\,KB & 47.1\,KB \\
4 & MFE (FFT 32) & 2D CNN & 99.6\% & 167\,ms & 30.2\,KB & 52.2\,KB \\
5 & MFE + clusters & 2D + Anom. & 99.7\% & 167\,ms & 30.2\,KB & 52.2\,KB \\
\bottomrule
\end{tabular}
\end{table}

\begin{table}[t]
\centering
\caption{Closed-split baselines (curated held-out audio).}
\label{tab:closed}
\small
\begin{tabular}{lcccc}
\toprule
\textbf{Model} & \textbf{Acc.} & \textbf{Prec.} & \textbf{Rec.} & \textbf{F1} \\
\midrule
SVM (RBF, MFCC) & 0.953 & 0.99 & 0.91 & 0.95 \\
Random Forest (MFCC) & 0.99 & 0.99 & 0.99 & 0.99 \\
1D CNN (MFCC) & 0.96 & 0.96 & 0.95 & 0.96 \\
2D CNN (MFE) & 0.98 & 0.98 & 0.97 & 0.98 \\
\bottomrule
\end{tabular}
\end{table}

\begin{table*}[t]
\centering
\caption{Architecture specifications and INT8 edge footprint.}
\label{tab:arch}
\small
\begin{tabular}{clcccc}
\toprule
\textbf{\#} & \textbf{Architecture} & \textbf{Input} & \textbf{Tensor} & \textbf{Params} & \textbf{INT8 Size} \\
\midrule
1 & Baseline 1D CNN & Raw waveform & $(1,16537,1)$ & 92{,}513 & 112.5\,KB \\
2 & Baseline 2D CNN & Log-Mel & $(1,64,130,1)$ & 249{,}985 & 259.9\,KB \\
3 & Enhanced 1D CNN & Raw, dual-head & $(1,16537,1)$ & 96{,}674 & 118.5\,KB \\
4 & Enhanced 2D CNN & Log-Mel & $(1,64,33,1)$ & 110{,}594 & 123.6\,KB \\
5 & Robust CRNN-PCEN & PCEN + unrolled GRU & $(1,64,130,1)$ & 397{,}842 & 478.9\,KB \\
\bottomrule
\end{tabular}
\end{table*}

\begin{table*}[t]
\centering
\caption{300-clip in-the-wild benchmark at $\tau = 0.50$ (sliding vs.\ peak-centered).}
\label{tab:bench}
\small
\begin{tabular}{lcccccc}
\toprule
\textbf{Model} & \multicolumn{3}{c}{\textbf{Sliding window}} & \multicolumn{3}{c}{\textbf{Peak-centered}} \\
\cmidrule(lr){2-4}\cmidrule(lr){5-7}
 & Recall & Amb.\ Spec. & Imp.\ Rej. & Recall & Amb.\ Spec. & Imp.\ Rej. \\
\midrule
Baseline 1D CNN & 45.0\% & 85.0\% & 72.0\% & 30.0\% & 95.0\% & 84.0\% \\
Baseline 2D CNN & 66.0\% & 88.0\% & 59.0\% & 45.0\% & 94.0\% & 77.0\% \\
Enhanced 1D CNN & \textit{100\% (deg.)} & \textit{0.0\%} & \textit{0.0\%} & \textit{100\% (deg.)} & \textit{0.0\%} & \textit{0.0\%} \\
Enhanced 2D CNN & \textbf{67.0\%} & 82.0\% & 66.0\% & 38.0\% & \textbf{99.0\%} & \textbf{89.0\%} \\
Robust CRNN-PCEN & 57.0\% & 86.0\% & 61.0\% & 33.0\% & 97.0\% & 85.0\% \\
\bottomrule
\end{tabular}
\end{table*}

\begin{table}[t]
\centering
\caption{CRNN-PCEN domain robustness (clean vs.\ synthetic mic distortion, sliding).}
\label{tab:domain}
\small
\begin{tabular}{cccccc}
\toprule
$\tau$ & \textbf{Clean Rec.} & \textbf{Dist.\ Rec.} & $\Delta$ & \textbf{Cl.\ Amb.} & \textbf{Dist.\ Amb.} \\
\midrule
0.10 & 69.0\% & 69.0\% & 0.0\% & 75.0\% & 74.0\% \\
0.30 & 62.0\% & 66.0\% & +4.0\% & 78.0\% & 77.0\% \\
0.50 & 61.0\% & 61.0\% & 0.0\% & 80.0\% & 77.0\% \\
0.70 & 59.0\% & 59.0\% & 0.0\% & 84.0\% & 81.0\% \\
\bottomrule
\end{tabular}
\end{table}

\begin{table}[t]
\centering
\caption{Multi-model ensemble 3$\times$3 confusion matrix (300 clips).}
\label{tab:ensemble}
\small
\begin{tabular}{lcccc}
\toprule
\textbf{Ground truth} & \textbf{Real} & \textbf{Fake} & \textbf{Not} & \textbf{Acc.} \\
\midrule
Real gunshot & \textbf{77} & 15 & 8 & 77.0\% \\
Hard imposter & 44 & \textbf{46} & 10 & 46.0\% \\
Ambient & 20 & 38 & \textbf{42} & 42.0\% \\
\bottomrule
\end{tabular}
\end{table}

\begin{table}[t]
\centering
\caption{Latency profiling (100 iterations, INT8).}
\label{tab:latency}
\small
\begin{tabular}{lccc}
\toprule
\textbf{Model} & \textbf{Mean} & \textbf{P95} & \textbf{Hop budget} \\
\midrule
Baseline 1D CNN & 2.29\,ms & 1.98\,ms & 1.2\% \\
Baseline 2D CNN & 2.15\,ms & 1.69\,ms & 1.2\% \\
Enhanced 1D CNN & 2.69\,ms & 6.48\,ms & 1.4\% \\
Enhanced 2D CNN & \textbf{0.09\,ms} & \textbf{0.10\,ms} & \textbf{0.05\%} \\
Robust CRNN-PCEN & 11.86\,ms & 14.50\,ms & 6.3\% \\
\bottomrule
\end{tabular}
\end{table}

\noindent\textbf{Findings.} Sliding mode buys $+15$ to $+29$ percentage points of recall at the cost of lower
specificity; peak mode buys up to $99\%$ specificity. This recall--specificity trade-off motivates the
ensemble. Under synthetic microphone distortion (bandpass 150--8000\,Hz, clipping $\pm0.60$--$0.95$, noise
15--30\,dB SNR, gain $0.5\times$--$1.5\times$) the CRNN-PCEN showed \textbf{zero} recall degradation and
\emph{improved} imposter rejection ($+4$--$7\%$), confirming PCEN as an intrinsic domain-shift normalizer. The
ensemble reached \textbf{77\%} recall --- for example, file \texttt{127724.wav} scored only $21.1\%$ on the
Enhanced 2D CNN but $99.3\%$ on the CRNN. The Enhanced 1D CNN is a cautionary case: after INT8 quantization
its raw-waveform dual-head logits saturated, giving $100\%$ recall with $0\%$ specificity. Host-vs-VM recall
differed by only $\pm4$--$6\%$ (INT8 rounding between XNNPACK and ARM NEON kernels). All five architectures
sit far below the $187.5$\,ms deadline, so zero buffer drops are guaranteed.

% ==============================================================================
\section{Discussion}
\label{sec:discussion}
% ==============================================================================
\subsection{The Honesty of the Journey}
Whether all this actually ``worked'' is still, in a way, a kind of mystery: a model can report
$100{:}100{:}100$ on test and validation while it only sits on a 1:1 set of gunshot and non-gunshot for the
next purpose. When the Raspberry Pi did not work well, we kept all those results, and what we now present is
the final one of all those tried steps.

\subsection{What Is Genuinely Unique}
\label{sec:unique}
Everybody does CNN, and everybody has PCEN --- we took PCEN from someone as well. What is unique in our
open-standard stack is how we combine it: a PCEN bioacoustic front-end for a gunshot detector, a measured
proof of domain-shift invariance, and a Voronoi-based spatial consensus layer at the sink that uses only open
standards (no proprietary OS), together with the correct layer depth reported in the literature (roughly
14--16 layers). This is the section for ``some uniqueness''; from here we move to the next sub-parts.

\subsection{The Single-Microphone Physical Boundary}
Single-microphone sensing has an inviolable limit: \textbf{chemical aerial mortar fireworks}. Certain
commercial fireworks produce combustion rise-times and reverberant decay tails that are acoustically
identical to a small-caliber muzzle blast on one capsule. Separating them cannot be done by audio
classification alone --- it needs spatial multi-microphone arrays measuring propagation velocity and
angle-of-arrival (TDOA).

% ==============================================================================
\section{Edge-Native WSN Deployment Architecture}
\label{sec:wsn}
% ==============================================================================
\subsection{The Sensor Node}
Each node is an \textbf{Arduino Nano 33 BLE Sense} (Nordic nRF52840, ARM Cortex-M4F @ 64\,MHz, 1\,MB Flash,
256\,KB SRAM) with an on-board digital PDM microphone; the quantized INT8 classifier runs entirely on-node.

\subsection{Two-Tier Sleep Scheduling}
Continuous full-power classification would drain a battery in 24--48\,h, so nodes use a two-tier duty cycle:
\textbf{Tier 1} deep sleep ($\sim 5\,\mu$A) with a low-power threshold/comparator ($< 75$\,dB SPL $\Rightarrow$
sleep); on an energy interrupt, \textbf{Tier 2} wakes at 64\,MHz, buffers a 250--500\,ms DMA window, runs the
INT8 model (1--167\,ms), decides, and returns to sleep.

\subsection{The 1-Bit Binary Alert}
Raw audio never leaves the node: streaming 16\,kHz/16-bit audio would need $\sim 256$\,kbps continuously per
node. Instead the output is effectively \textbf{binary}: if it is a gunshot the node transmits, and if it is
not, it stays silent. On confidence above threshold it broadcasts
\begin{equation}
    \text{Payload} = \{\text{Node\_ID (1\,B)}, \text{Timestamp (4\,B)}, \text{Detection (1\,bit)}, \text{Confidence (1\,B)}\},
\end{equation}
and otherwise transmits nothing. Radio on-time is $<1.5$\,ms, driving RF energy to near zero
($<0.01$\,mWh/event) and giving multi-year battery life on a coin cell. Emitting only a $1$ for ``gunshot''
and a $0$ (silence) otherwise is, in effect, a \textbf{binary neural-network output} --- the most energy- and
privacy-friendly signal in the system.

\subsection{Sink Aggregation (Raspberry Pi Hub)}
A lone trigger is only a candidate. The hub enforces spatial consensus and localisation:
\textbf{Voronoi} clustering of nodes into cells; a \textbf{coincidence window} validating an incident only if
$K \ge 3$ neighbouring nodes report detection within
\begin{equation}
    \Delta t \le \frac{d_{\max}}{c_{\text{sound}}} \approx \frac{34\ \text{m}}{343\ \text{m/s}} \approx 100\ \text{ms},
\end{equation}
and \textbf{hyperbolic TDOA localisation}:
\begin{equation}
    \sqrt{(x-x_i)^2+(y-y_i)^2} - \sqrt{(x-x_j)^2+(y-y_j)^2} = c\,(t_i-t_j).
\end{equation}
This discards isolated impacts and near-field claps and gives the network its fault-tolerant (self-healing)
property: one noisy or failed node cannot raise a false alarm, and the mesh re-routes around it.

\subsection{Dashboard and Privacy}
Verified, localised events are pushed to a live \textbf{dashboard} (map $+$ SMS push for field personnel).
Because conversational audio is processed strictly inside volatile SRAM and immediately overwritten, the
system is architecturally immune to eavesdropping --- the privacy guarantee is mathematical, not policy.

% ==============================================================================
\section{Conclusion and Future Work}
\label{sec:conclusion}
% ==============================================================================
This project shows, end to end, that an acoustic gunshot detection system can be built on inexpensive,
widely available hardware without any cloud dependency. Starting from a raw UrbanSound-class corpus we built
a reproducible Data-Cleaner pipeline, learned the hard way why 250\,ms trimming and MFCC temporal averaging
fail, migrated to an MFE spectrogram pipeline, ran a full Edge Impulse TinyML study, advanced through the CNN
family (1D $\rightarrow$ 2D $\rightarrow$ 3D $\rightarrow$ enhanced dual-head), and finally adopted a
bioacoustic PCEN $+$ Bi-GRU model that we quantized and unrolled into a native INT8 LiteRT engine. On an
unseen 300-clip in-the-wild benchmark the Enhanced 2D CNN delivered sub-millisecond inference with up to
$99\%$ ambient specificity, the CRNN-PCEN proved zero recall loss under microphone domain shift, and a
multi-model ensemble reached $77\%$ gunshot recall with physics-based imposter rejection. We closed by trying
it finally on the (virtual) Raspberry Pi, and it did make much of this perfect.

The remaining work runs along two tracks. The first is hardware: our next step is the course of work related
to the \textbf{hardware}, moving from the virtual Raspberry Pi to real boards and finally onto the
\textbf{Arduino Nano 33 BLE} for on-device inference --- our final goal is to deploy, and that is it. The
second is to make the real-world IoT system as optimal as it can be: deeper \textbf{sleep scheduling}, longer
\textbf{battery life}, and \textbf{self-healing} behaviour across the mesh (redundant nodes, consensus, and
graceful recovery), with the whole network built on standard \textbf{WSN/IoT} concepts and surfaced through a
live \textbf{dashboard}. The most powerful idea of all is to keep the over-the-air result to a \textbf{single
bit} --- transmit $1$ only when a gunshot is fired, and stay silent for $0$ --- which is the binary inference
that saves the most energy. Keeping only the best and most relevant of our many result files, this is the
shape of the six-page research paper we intend to publish.

% ==============================================================================
\begin{thebibliography}{99}
\bibitem{dib2023} Data in Brief, ``A multi-firearm, multi-orientation audio dataset of gunshots,'' \textit{Data in Brief}, vol. 48, p. 109091, 2023. DOI: 10.1016/j.dib.2023.109091.
\bibitem{lostanlen2019} V. Lostanlen, K. Cramer, S. Farnsworth, F. Briggs, and J. Bello, ``Per-Channel Energy Normalization for Bioacoustic Sound Event Detection,'' \textit{IEEE/ACM Trans. Audio, Speech, Lang. Process.}, vol. 27, no. 12, pp. 2178--2190, 2019.
\bibitem{hershey2017} S. Hershey et al., ``CNN Architectures for Large-Scale Audio Classification,'' in \textit{Proc. IEEE ICASSP}, 2017, pp. 131--135.
\bibitem{zhang2018} H. Zhang, M. Cisse, Y. N. Dauphin, and D. Lopez-Paz, ``mixup: Beyond Empirical Risk Minimization,'' in \textit{Proc. ICLR}, 2018.
\bibitem{park2019} D. S. Park et al., ``SpecAugment: A Simple Data Augmentation Method for Automatic Speech Recognition,'' in \textit{Proc. Interspeech}, 2019, pp. 2613--2617.
\bibitem{edgeimpulse2024} Edge Impulse, ``TinyML Acoustic Anomaly Detection on Microcontrollers,'' Tech. Rep., 2024.
\bibitem{magee2019} G. Magee, ``Gunshot Detection System on Raspberry Pi,'' IEEE Open Source Acoustic Repository, 2019.
\bibitem{saha2025} S. Saha et al., ``Automated Gunshot Detection in Forest Environments Using Spectrogram CNNs,'' \textit{Applied Acoustics}, 2025.
\end{thebibliography}

\end{document}
